\documentclass[11pt]{article}
\usepackage[margin=1in]{geometry}
\usepackage{amsmath,amssymb,bm,booktabs,array,microtype}
\usepackage[hidelinks]{hyperref}
\usepackage{xcolor}
\newcommand{\tr}{\operatorname{tr}}
\newcommand{\diag}{\operatorname{diag}}
\newcommand{\E}{\mathbb E}
\newcommand{\R}{\mathbb R}
\newcommand{\bI}{\bm I}
\newcommand{\bP}{\bm P}
\newcommand{\bK}{\bm K}
\newcommand{\by}{\bm y}
\newcommand{\be}{\bm e}
\newcommand{\bX}{\bm X}
\newcommand{\bW}{\bm W}
\newcommand{\bC}{\bm C}
\newcommand{\bD}{\bm D}
\newcommand{\bone}{\bm 1}

\title{A reference-cohort summary estimator for genetic and gene--environment variance components}
\author{SUMMIT technical report}
\date{14 August 2026}

\begin{document}
\maketitle

\begin{abstract}
GENIE estimates additive genetic (G), gene--environment (G$\times$E),
noise--environment (N$\times$E), and residual variance components from one
individual-level cohort.  Rewriting GENIE's right-hand side in terms of
marginal GWAS and GWIS scores is exact when the kernel traces and the scores
come from the same individuals.  It is not, by itself, a reusable reference
method.  This report derives the additional moments needed when G and
G$\times$E kernel quantities are estimated in an independent cohort containing
all individuals with measured environment, while trait-specific marginal
scores may use smaller and different samples.  The key is to transfer the
same-individual and different-individual parts of each genetic kernel-product
trace separately.  For G$\times$E, the same-individual term is a non-universal
fourth moment and cannot be replaced by the additive-only constant.  It can be
estimated during the existing randomized, in-memory reference pass with an
order-two probe U-statistic and $O(N_RK)$ additional memory, independent of the
number of probes.  Study-specific N$\times$E moments remain exact low-dimensional
summaries.  We also audit covariate handling in the released GENIE code and
distinguish two coherent feature definitions from the incoherent mixture of
raw projected kernels and hard-coded traces.
\end{abstract}

\section{Notation and generative assumptions}

We follow the notation of Kerin and Marchini~\cite{genie}.  Cohort $S$ is the
trait study and cohort $R$ is the reusable reference.  They share an ordered
set of $M$ variants and the same annotation weights $A_{jk}\geq0$,
$k=1,\ldots,K$, with $M_k=\sum_j A_{jk}$.  Their individuals need not be the
same.  Let $\bC_S$ contain an intercept, the environment, and the other fixed
effects used by the trait analysis.  Define
\begin{equation}
 \bP_S=\bI-\bC_S(\bC_S^\top\bC_S)^+\bC_S^\top,
 \qquad r_S=\tr(\bP_S)=N_S-\operatorname{rank}(\bC_S),
 \label{eq:projector}
\end{equation}
where $(\cdot)^+$ is the Moore--Penrose inverse.  The same definitions apply
to $R$.

Let $\bm G_S$ be a centered, pre-scaled genotype matrix and let $\be_S$ be a
centered environment measured in standard-deviation units.  The default
one-environment GENIE generative model, before fixed-effect projection, is
\begin{equation}
 \widetilde{\by}_S=\bm G_S\widetilde{\bm\beta}
  +\diag(\be_S)\bm G_S\widetilde{\bm\alpha}
  +\diag(\be_S)\bm\delta+\bC_S\bm\gamma+\bm\varepsilon.
 \label{eq:rawmodel}
\end{equation}
The published framework allows several environments; the present derivation
applies to one environment at a time.  GENIE uses the naturally scaled
projected columns $\bP_S\bm G_S$ and
$\bP_S\diag(\be_S)\bm G_S$.  SUMMIT's default instead restandardizes those
columns after projection:
\begin{align}
 \bm U_S&=\bP_S\bm G_S,&
 \bX_{S,:,j}&=\bm U_{S,:,j}
       \sqrt{\frac{r_S}{\bm U_{S,:,j}^\top\bm U_{S,:,j}}},\nonumber\\
 \bm V_S&=\bP_S\diag(\be_S)\bm G_S,&
 \bW_{S,:,j}&=\bm V_{S,:,j}
       \sqrt{\frac{r_S}{\bm V_{S,:,j}^\top\bm V_{S,:,j}}}.
 \label{eq:features}
\end{align}
Thus every retained column of $\bX_S$ and $\bW_S$ has squared norm $r_S$.
This is the G$\times$E analogue of SUMMIT's covariate-residualized,
restandardized additive feature definition~\cite{summit}.  The two choices
define different random-effect scalings, but the method-of-moments algebra
below is the same when all traces and scores use one choice coherently.  The
external implementation uses the standardized choice.  Define
\begin{align}
 \bK^G_{Sk}&=\frac{\bX_S\diag(\bm A_{:k})\bX_S^\top}{M_k},&
 \bK^I_{Sk}&=\frac{\bW_S\diag(\bm A_{:k})\bW_S^\top}{M_k},\nonumber\\
 \bD_S&=\bP_S\diag(\be_S^2)\bP_S,&
 \bK^0_S&=\bP_S.
 \label{eq:kernels}
\end{align}
The superscript $I$ denotes interaction and avoids confusing an environment
vector with an expectation.

The residualized phenotype is generated under the second-moment model
\begin{equation}
 \by_S=\sum_{k=1}^K\bX_S\bm\beta_k
       +\sum_{k=1}^K\bW_S\bm\alpha_k
       +\bP_S\diag(\be_S)\bm\delta+\bP_S\bm\varepsilon,
 \label{eq:model}
\end{equation}
with mutually uncorrelated, mean-zero effects
\begin{align*}
 \operatorname{Cov}(\bm\beta_k)&=\frac{\sigma^2_{Gk}}{M_k}
       \diag(\bm A_{:k}),&
 \operatorname{Cov}(\bm\alpha_k)&=\frac{\sigma^2_{Ik}}{M_k}
       \diag(\bm A_{:k}),\\
 \operatorname{Cov}(\bm\delta)&=\sigma_D^2\bI,&
 \operatorname{Cov}(\bm\varepsilon)&=\sigma_0^2\bI.
\end{align*}
For overlapping annotations this is understood as a linear covariance model;
the coefficients need not each be nonnegative.  Conditional on the design,
\begin{equation}
 \bm\Sigma_S=\operatorname{Cov}(\by_S\mid\bX_S,\bW_S,\be_S,\bC_S)
 =\sum_k\sigma^2_{Gk}\bK^G_{Sk}
 +\sum_k\sigma^2_{Ik}\bK^I_{Sk}
 +\sigma_D^2\bD_S+\sigma_0^2\bP_S.
 \label{eq:covariance}
\end{equation}
No Gaussian assumption on the effects is required for the moment equations.

\section{GENIE as an in-sample method-of-moments estimator}

Index the $2K+2$ kernels in Eq.~\eqref{eq:covariance} by
$\mathcal K_S=(\bK^G_{S1},\ldots,\bK^G_{SK},
\bK^I_{S1},\ldots,\bK^I_{SK},\bD_S,\bP_S)$ and collect their coefficients in
$\bm\theta$.  GENIE minimizes the Frobenius loss
\begin{equation}
 \left\|\by_S\by_S^\top-\sum_a\theta_a\mathcal K_{Sa}\right\|_F^2.
\end{equation}
Differentiation gives the normal equations
\begin{equation}
 \bm T_S\widehat{\bm\theta}=\bm q_S,
 \qquad (T_S)_{ab}=\tr(\mathcal K_{Sa}\mathcal K_{Sb}),
 \qquad (q_S)_a=\by_S^\top\mathcal K_{Sa}\by_S.
 \label{eq:normal}
\end{equation}
This is precisely Supplemental Note~S1 of GENIE after writing the projected
kernels explicitly~\cite{geniesupp}.  GENIE estimates large traces by
Hutchinson random projection: for independent isotropic probes $\bm z_v$,
\begin{equation}
 \tr(\mathcal K_a\mathcal K_b)
 =\E_v[(\mathcal K_a\bm z_v)^\top(\mathcal K_b\bm z_v)].
\end{equation}

For a genetic family $F\in\{G,I\}$, define the marginal HE-scale score
\begin{equation}
 s^F_{Sj}=\frac{\bm F_{S,:,j}^\top\by_S}{\sqrt{r_S}}.
 \label{eq:score}
\end{equation}
It follows exactly that
\begin{equation}
 \by_S^\top\bK^F_{Sk}\by_S
 =\frac{r_S}{M_k}\sum_j A_{jk}(s^F_{Sj})^2.
 \label{eq:rhs}
\end{equation}
The required statistics are marginal cross-products after the common
fixed-effect projection.  A usual conditional interaction Wald statistic is
not interchangeable with Eq.~\eqref{eq:score} unless its numerator and scale
are transformed back to this marginal score.

Let
\begin{equation}
 L^{FH}_{S,jk}=\sum_m A_{mk}
 \left(\frac{\bm F_{S,:,j}^\top\bm H_{S,:,m}}{r_S}\right)^2.
 \label{eq:ldpanel}
\end{equation}
Then the genetic block of Eq.~\eqref{eq:normal} is also exact:
\begin{equation}
 \tr(\bK^F_{Sk}\bK^H_{S\ell})
 =\frac{r_S^2}{M_kM_\ell}\sum_j A_{jk}L^{FH}_{S,j\ell}.
 \label{eq:exacttrace}
\end{equation}
Equations~\eqref{eq:rhs} and \eqref{eq:exacttrace} explain the existing
matched-cohort summary implementation.  They do \emph{not} justify inserting
an independently estimated $L_R$ into Eq.~\eqref{eq:exacttrace}: the
finite-cohort trace has both same-person fourth moments and different-person
LD moments, which scale differently with cohort size.

\section{A genuinely reusable reference estimator}

\subsection{Population transport assumptions}

The reference construction requires the following assumptions.
\begin{enumerate}
 \item Cohorts $R$ and $S$ are independent draws from the same target
 population, with the same variant/allele and annotation definitions.
 \item After applying the declared fixed-effect convention and per-cohort
 restandardization, the joint second and fourth moments of $(X,W,E)$ are
 transportable.  Fixed-effect dimension is small relative to cohort size, so
 the projection's dependence and leverage perturb the exchangeable-individual
 moments only at lower order.  Large leverage differences or
 population/LD/environment shift violate this assumption.
 \item Trait missingness may change $N_S$ but is ignorable conditional on the
 modeled variables.  Genotypes, environment, and covariates used to form the
 marginal scores obey the same feature convention as the reference.
\end{enumerate}
These are reference-panel assumptions, not algebraic identities.  If the
joint G--E distribution differs, the estimator converges to a reference-weighted
pseudo-parameter, just as additive SUMMIT with mismatched LD has target
$(\bm A_R)^{-1}\bm A_S\bm h^2$~\cite{summit}.

\subsection{Same-person and different-person trace moments}

For any genetic kernels $a=(F,k)$ and $b=(H,\ell)$, write
\begin{align}
 D^C_{ab}&=\sum_{i=1}^{N_C}K^C_a(i,i)K^C_b(i,i),\nonumber\\
 O^C_{ab}&=\sum_{i\ne i'}K^C_a(i,i')K^C_b(i,i'),\qquad
 T^C_{ab}=D^C_{ab}+O^C_{ab}.
 \label{eq:do}
\end{align}
The split is on the individual axis, not an arbitrary orthonormal basis for
the residual subspace.  Residual rank $r_C$ controls feature and score
normalization, but it is not the number of exchangeable people.  Define
population moments
\begin{equation}
 d_{ab}=\E[K_a(i,i)K_b(i,i)],\qquad
 o_{ab}=\E[K_a(i,i')K_b(i,i')],\quad i\ne i'.
\end{equation}
Then
\begin{equation}
 \E[T^S_{ab}]=N_Sd_{ab}+N_S(N_S-1)o_{ab}.
 \label{eq:populationtrace}
\end{equation}
The reference estimators are
\begin{equation}
 \widehat d_{ab}=\frac{\widehat D^R_{ab}}{N_R},\qquad
 \widehat o_{ab}=\frac{\widehat T^R_{ab}-\widehat D^R_{ab}}
 {N_R(N_R-1)},
\end{equation}
and therefore
\begin{equation}
 \boxed{\widehat T^{S\mid R}_{ab}
 =\frac{N_S}{N_R}\widehat D^R_{ab}
 +\frac{N_S(N_S-1)}{N_R(N_R-1)}
  (\widehat T^R_{ab}-\widehat D^R_{ab}).}
 \label{eq:transfer}
\end{equation}
When $N_S=N_R$, Eq.~\eqref{eq:transfer} returns the original reference trace
exactly.  For additive standardized genotypes, $D^R$ is close to the familiar
$N_R$-order baseline.  For G$\times$E, it contains quantities such as
$\E[X^2E^2X'^2E'^2]$ and is not universal.  This is the precise reason that
blind squared-rank scaling, or importing the additive leading term, fails.
There is no assumption-free way to recover a study's conditional trace from
an independent cohort.  Equation~\eqref{eq:transfer} is the reference-panel
estimator under the population-transport assumptions above; without them the
target is necessarily reference weighted.

\subsection{Estimating the fourth-moment term without a second genotype pass}

For feature family $a=(F,k)$ and probe $v$, the reference trace pass already
forms
\begin{equation}
 S_{aiv}=\sum_jF_{R,ij}\sqrt{A_{jk}}z_{jv},
 \qquad \E_z[S_{aiv}^2]=M_kK^R_a(i,i).
\end{equation}
For $B\ge2$, an unbiased order-two probe U-statistic is
\begin{equation}
 \widehat D^R_{ab}=
 \frac{1}{B(B-1)M_kM_\ell}
 \sum_i\left[
  \left(\sum_vS_{aiv}^2\right)\left(\sum_uS_{biu}^2\right)
  -\sum_vS_{aiv}^2S_{biv}^2
 \right].
 \label{eq:ustat}
\end{equation}
The subtraction excludes identical probes; different probe indices are
independent even when annotations overlap and G/G$\times$E sketches share the
same underlying probe.  An implementation needs an $N_R\times2K$ matrix of
running square sums and one $2K\times2K$ accumulator.  Temporary conversion is
chunked over probes.  Memory is $16N_RK+O(K^2)$ bytes in float64 and is
independent of $B$; at $N_R=291{,}273,K=1$ it is about 4.4 MiB.  No source
sketch, block sketch, or cache is written to disk.

\subsection{N$\times$E and residual moments}

N$\times$E is study specific and cheap; it should not be borrowed from $R$.
Let $\bm Q_S$ be an orthonormal basis for the fixed effects and
$\bm d_S=\be_S^2$.  With $\bD_S=\bP_S\diag(\bm d_S)\bP_S$,
\begin{align}
 \tr(\bD_S)
 &=\sum_i d_{Si}-\tr(\bm Q_S^\top\diag(\bm d_S)\bm Q_S),\label{eq:trd}\\
 \tr(\bD_S^2)
 &=\sum_i d_{Si}^2
 -2\tr(\bm Q_S^\top\diag(\bm d_S^2)\bm Q_S)
 +\tr\{(\bm Q_S^\top\diag(\bm d_S)\bm Q_S)^2\}.
 \label{eq:trd2}
\end{align}
The phenotype supplies
\begin{equation}
 q_D=\by_S^\top\bD_S\by_S=\|\be_S\odot\by_S\|^2,
 \qquad q_0=\by_S^\top\by_S=r_S,
\end{equation}
on SUMMIT's normalized residual scale.  Cross-traces with genetic kernels are
also exact study summaries.  The same pass that computes marginal GWAS/GWIS
scores forms
\begin{equation}
 d^F_{Sj}=\frac{\bm F_{S,:,j}^{\top}\diag(\be_S^2)
                    \bm F_{S,:,j}}{r_S},\qquad
 \tr(\bK^F_{Sk}\bD_S)=\frac{r_S}{M_k}\sum_j A_{jk}d^F_{Sj},
 \label{eq:kd}
\end{equation}
because projected feature columns satisfy
$\bP_S\bm F_{S,:,j}=\bm F_{S,:,j}$.  This costs no extra genotype decode and
avoids transporting an environment-sensitive fourth moment.  Because
Eq.~\eqref{eq:features} makes
$\tr(\bK^S_a)=r_S$, genetic--residual cross-traces are exact.

Equations~\eqref{eq:rhs}, \eqref{eq:transfer}, \eqref{eq:trd}--\eqref{eq:kd}
fill every entry of the study normal equations.  Solving them estimates the
four classes $(G,G\times E,N\times E,\text{residual})$ using genotype/environment
moments from all eligible reference individuals and trait-specific marginal
scores from the study cohort.

\subsection{The summary-data boundary}

The reusable reference object contains the ordered variant/allele and
annotation definitions, $N_R,r_R$, the randomized genetic trace matrix
$\widehat{\bm T}_R$, and the compact same-individual matrix
$\widehat{\bm D}_R$.  These quantities are phenotype independent and are
computed once in every individual with an observed environment and reference
covariates.  For each trait, the required study object is
\begin{equation}
 \mathcal S_y=\{N_S,r_S,\bm s^G_S,\bm s^I_S,
 q_D,q_0,\tr(\bD_S),\tr(\bD_S^2),
 \tr(\bK^G_{Sk}\bD_S),\tr(\bK^I_{Sk}\bD_S)\}_{k=1}^K.
 \label{eq:summaryinterface}
\end{equation}
The two score vectors are marginal GWAS/GWIS cross-products.  The remaining
terms are scalars or $O(K)$ design summaries.  Equation~\eqref{eq:kd} lets the
genetic--N$\times$E terms be accumulated in the same streamed genotype pass
that produces GWAS/GWIS, with no randomized trace calculation and no second
decode.  Once $\mathcal S_y$ is written, fitting uses no individual-level
phenotype, genotype, or environment.  Thus trait missingness affects the cheap
trait summary but does not force reconstruction of the all-E G$\times$E
reference.

\subsection{Uncertainty and identifiability}

There are three uncertainty sources: trait-score sampling, finite reference
sampling, and randomized-probe error.  A SNP-block jackknife deletes the
corresponding GWAS/GWIS score rows and reference LD-score rows.  The default
block-local approximation assumes cross-block LD is negligible and reuses the
full-reference $\widehat D^R$ in each replicate.  This omits a change of order
the deleted SNP fraction in the diagonal term, so it is an explicit
approximation whose stability should be checked when blocks are few, large, or
annotation-imbalanced.  An exact deletion would require block-specific
cross-sketch moments and is intentionally not the production default.

Increasing $B$ reduces only probe error.  It cannot repair reference mismatch
or an unidentified model.  In particular, if $E^2$ is constant after
projection, $\bD_S$ is proportional to $\bP_S$ and N$\times$E is aliased with
residual noise.  Collinear annotation kernels cause analogous rank failure.

\section{Validation}

Three independent checks were used.
\begin{enumerate}
 \item A dense oracle forms all individual-level kernels and verifies that the
 matched summary equations reproduce every GENIE matrix, right-hand side, and
 trace entry to floating-point tolerance (maximum tested discrepancy below
 $2\times10^{-13}$).
 \item Setting $N_S=N_R$ in the implemented transfer reproduces the complete
 reference genetic block exactly, irrespective of how $T_R$ is split into
 $D_R$ and $O_R$.
 \item In a separate continuous-$E$ Monte Carlo experiment, an independent
 $N_R=2500$ reference and 60 independent $N_S=350$ studies used $M=60$
 correlated Gaussian markers (AR(1) correlation $0.35$), with intercept and
 $E$ projected out.  Median relative errors of Eq.~\eqref{eq:transfer} were
 2.17\% (G--G), 3.14\% (G--G$\times$E), and 2.23\%
 (G$\times$E--G$\times$E).  Blind squared-rank scaling gave 9.27\%, 86.46\%,
 and 23.51\%, respectively.  Genetic--N$\times$E cross-traces were computed
 exactly in each study rather than transported.  These are finite-reference
 errors, not claimed universal operating characteristics.
\end{enumerate}
An end-to-end integration test also constructs one reference, removes a trait
value to create a different study cohort, writes the study GWAS/GWIS moments,
and fits the transferred equations.

\section{Covariates and restandardization in released GENIE code}

The equations in the GENIE paper are coherent: with
$\bm V=\bI-\bC(\bC^\top\bC)^{-1}\bC^\top$, Supplemental Note~S1 uses
$\tr(K_kVK_\ell V)$, $\tr(VK_k)$, $\by^\top VK_kV\by$, and residual dimension
$N-C$~\cite{geniesupp}.  The concerns below apply to the released code at
commit \texttt{22bee9c}, not to that formal derivation.

\paragraph{Ordinary inverse and column count.}
The implementation computes \texttt{Q = WtW.inverse()} and sets the residual
dimension to \texttt{Nindv\_mask - Ncov}~\cite{geniecodecov}.  This is correct
only if the retained covariate columns are exactly full rank and numerically
well conditioned.  Environment, intercept, duplicated indicators, or a
constant covariate can make that false.  A rank-revealing QR/SVD and
$N-\operatorname{rank}(C)$ are required for the projector in
Eq.~\eqref{eq:projector}.

\paragraph{Scaling before the analysis mask.}
The genotype code computes allele means over \texttt{Nindv} and applies the
HWE scale $\{2p(1-p)\}^{-1/2}$ before the phenotype/environment mask and
fixed-effect projection~\cite{geniecodescale}.  Consequently a retained,
projected column generally does not have squared norm $N_{\rm mask}$ or
$N_{\rm mask}-\operatorname{rank}(C)$.  This is especially visible when trait
missingness creates a much smaller analysis sample.

\paragraph{The important distinction: two coherent estimands and one
incoherent hybrid.}
There are two valid choices.
\begin{enumerate}
 \item \emph{Raw projected kernel:} use $U=PX$ and $W=P(E\odot X)$ without
 post-projection scaling.  This defines a valid model, but every trace,
 including $\tr(K)$, must be computed from the realized columns and the random
 effect prior is tied to the pre-projection genotype scale.
 \item \emph{Restandardized projected kernel:} rescale every projected column
 as in Eq.~\eqref{eq:features}.  Then all genetic kernel traces equal $r$ and
 summary moments have a stable, cohort-comparable per-residual-SD definition.
 This is SUMMIT's default.
\end{enumerate}
The released GENIE implementation initializes each additive genetic trace to
\texttt{Nindv\_mask}, while interaction/noise traces are estimated from the
random projections, and then applies a projection subtraction~\cite{geniecodetrace}.
That mixes the raw-kernel Gram and phenotype moments with a unit-norm additive
trace assumption that the preceding HWE scaling/masking does not guarantee.
The problem is not simply ``GENIE fails to restandardize'': raw projected
kernels are legitimate.  The problem is using their realized Gram/RHS while
hard-coding a trace from the alternative standardized-kernel convention.

\paragraph{Phenotype scaling.}
The default code standardizes the phenotype, projects covariates, and then
standardizes the projected phenotype again using denominator
$N_{\rm mask}-1$~\cite{geniecodepheno}, while the residual kernel dimension is
declared $N_{\rm mask}-N_{\rm cov}$.  Because the code retains the realized
$\by^\top\by$ in the normal equations, phenotype rescaling alone is not
necessarily a bias.  It does, however, mean that reported components are not
automatically on the $r$-normalized scale and should not be combined with
hard-coded $r$ traces without an explicit scale conversion.

These issues can be small when covariates are full rank, missingness is mild,
and HWE scaling nearly matches the retained sample variance.  They are not
algebraically controlled in the released implementation.  SUMMIT instead uses
a rank-revealing orthonormal basis, records the numerical rank, applies the
trait/reference mask before feature construction, and either restandardizes
projected columns or computes their actual traces consistently.

\section{Computational implications}

The new population statistic does not change the dominant complexity.  With
$B$ probes, $K$ annotations, $M$ variants, and streamed genotype blocks, the
G/G$\times$E reference remains linear in $N_R$, $M$, $B$, and $K$ for $K=1$
(and in the number of requested directional annotation products generally).
The added Eq.~\eqref{eq:ustat} work is $O(N_RBK^2)$ but has a very small
constant and no genotype decode or large matrix multiply; its persistent
memory is $O(N_RK)$.

On tabla, a sealed full $N_R=291{,}273$, $M=454{,}207$, $B=32$ EUR age
reference completed its numerical work in 1,503.6 seconds; end-to-end measured
elapsed time was 1,534.9 seconds (25.6 minutes), with 4.04 GiB peak RSS.  Its
feature, source, panel-preparation, and target phases took 560.0, 368.6, 0.20,
and 471.4 seconds.  The source tree, loaded native binary, and single OpenBLAS
0.3.34 runtime are hash-bound in the published reference manifest.  Relative
to the preceding internally threaded diagnostic, total elapsed time fell
20.9\%, target time fell 51.9\%, and repaired output columns fell from 384 to
3 without increasing memory.

A cohort-matched additive-only $B=32$ run used the same $N=291{,}273$,
$M=454{,}207$, covariates, probe distribution and seed, 32-core affinity, and
16-GiB sketch budget.  It completed in 568.4 seconds end to end (565.4 seconds
reported numerical runtime) with 1.39 GiB peak RSS.  The corresponding G$\times$E
ratios are therefore 2.70-fold for elapsed time and 2.90-fold for peak RSS.
The extra work is intrinsic to constructing the interaction features and
estimating the four directional G/G$\times$E kernel products instead of the one
additive product; it is not disk-cache I/O.  The comparison produced only a
4.97-MB additive LD-score file in addition to its small log and receipt.

A direct $B/32$ extrapolation is too pessimistic because the $B=32$ products
are skinny.  In a dense-oracle benchmark with $N=100{,}000$ and $M=1{,}000$,
increasing $B$ from 32 to 512 increased native source time from 0.319 to 1.267
seconds (3.97-fold) and the relevant 2$B$ target time from 0.458 to 2.396
seconds (5.23-fold), rather than 16-fold.  A subsequent direct $B=1024$ run
measured 2.026 seconds for the source product, 4.485 seconds for the 2$B$
target, and 8.453 seconds for the 4$B$ target.  Its maximum target error against
the independent dense float64 oracle was $2.538\times10^{-7}$, with no repaired
columns.  Scaling the sealed full-cohort $B=32$ phases by these measured
$B=1024/B=32$ ratios gives approximately 2.1 hours.  A conservative current
planning range is 2.2--2.8 hours to allow for repair and system-load variation.
The $B=1024$ component times are measured, but the full 454,207-variant time is
still a scaling estimate and should be replaced by an end-to-end measurement
before a large production campaign.

Probe tiling and block-width adaptation bound memory.  For one annotation, the
float32 $[XZ,WZ]$ source sketch costs $8N_RB$ bytes.  The native opaque
$[S,E\odot S]$ target panel costs $32N_RB$ bytes, while the protected
non-decoded target operand costs at most another $16N_RB$ bytes \emph{inside}
the configured native workspace.  At $N_R=291{,}273$ and $B=1024$, these terms
are about 2.22, 8.89, and 4.44 GiB.  The first two allocations plus the bounded
16-GiB native workspace (whose variant block shrinks as the protected operand
grows) imply an expected peak of roughly 25--30 GiB, not $32$ times the $B=32$
RSS.  The new population-transfer statistic itself adds only about 4.4 MiB for
$K=1$.

The numerical failures are not caused by extreme Rademacher probes: their
entries are exactly $\pm1$.  Nor is float32 the arithmetic root cause.  The
decoder, feature moments, projection, source/target GEMMs, checksums, and
repairs use float64; float32 controls only retained sketch buffers.  Switching
those buffers to float64 doubles their memory and reduces accumulation
roundoff, but cannot prevent a vendor GEMM from sparsely changing a double
input or output.  Dropping extreme LD-score values would be unsafe here: an
extreme Monte Carlo realization and a corrupt output are different events,
and clipping squared traces creates downward bias.

Before the environment repair, tabla loaded two BLAS runtimes in one process:
NumPy's pip wheel bundled OpenBLAS 0.3.30/Haswell while the native extension
used conda OpenBLAS 0.3.34/Zen.  The environment has now been changed to conda
NumPy 2.3.5, so NumPy and the extension both resolve to the same OpenBLAS
0.3.34 Zen library.  This removes a concrete source of process-wide thread and
runtime interference, but it did not remove the numerical fault.  The
full-cohort $B=32$ diagnostic triggered 384 repaired output columns.  More
decisively, a production-shaped dense test with $N=291{,}273$, $M=1{,}207$, and
$B=32$ intermittently observed a maximum target error of 58.15 from the
internally threaded vendor path, despite using binary64 arithmetic.  Repeating
the same test could also return only ordinary accumulation-order error, proving
that the failure is intermittent and shape/thread dependent.

SUMMIT therefore no longer invokes an internally threaded OpenBLAS GEMM for
these products.  It fixes OpenBLAS at one thread per call and partitions
disjoint output rows across SUMMIT-owned OpenMP workers.  This was faster, not
slower, in the production-shaped no-fault comparison (1.09 versus 1.33 seconds
for the 2$B$ target and 1.79 versus 2.28 seconds for 4$B$).  Across 20 repeated
2$B$/4$B$ calls, the selective ABFT layer repaired 257 output columns and every
final output remained within $2.283\times10^{-6}$ of the independent dense
float64 oracle.  Thus the vendor/concurrency problem is still observable, but
the returned result is protected without duplicating the full genotype
operand.  The ABFT path checks wide GEMMs, repairs only discrepant output
columns with a non-vendor cache-tiled kernel, and falls back to a deterministic
fresh decode only if the reconstructable genotype input changes.

\section{Practical contract}

The implemented external path is explicit.  A reference generated with at
least two probes records the compact population trace matrix in its manifest.
Trait scoring with \texttt{--gxe-population-}\allowbreak\texttt{reference} permits a different
retained sample, recomputes study feature standardization, and writes study
$N$, $r$, $q_D$, $q_0$, $\tr(D)$, $\tr(D^2)$, and exact genetic--$D$
cross-traces (plus delete-block versions when requested).  The fit requires the exact
reference-manifest hash, identical ordered variants/alleles and annotations,
the same named environment/covariate convention, and standardized kernels.
It fails closed for GENIE/raw kernels because their genetic traces are not
fixed by $r$.  Existing matched-cohort behavior remains the default.

\begin{thebibliography}{9}
\bibitem{genie}
Kerin M, Marchini J. Estimation of gene--environment interaction heritability
from large-scale biobanks. \emph{American Journal of Human Genetics} (2024).
\href{https://doi.org/10.1016/j.ajhg.2024.05.015}{doi:10.1016/j.ajhg.2024.05.015}.

\bibitem{geniesupp}
Kerin M, Marchini J. GENIE Supplemental Notes, especially Note S1.
\href{https://ars.els-cdn.com/content/image/1-s2.0-S0002929724001782-mmc1.pdf}{Publisher supplement}.

\bibitem{summit}
SUMMIT manuscript and Supplementary Information, local working version,
Methods ``Estimation of genome-wide LD scores'' and Supplementary Note
``Effect of reference-LD mismatch on SUMMIT estimates'' (accessed 13 August
2026).

\bibitem{geniecodecov}
GENIE source, commit \texttt{22bee9c}, covariate inverse and projection.
\href{https://github.com/sriramlab/GENIE/blob/22bee9cd53d61eabb2ffcec70acf3c5bb43ede01/src/ge_flexible.cpp#L2303-L2362}{\texttt{ge\_flexible.cpp}, lines 2303--2362}.

\bibitem{geniecodescale}
GENIE source, commit \texttt{22bee9c}, genotype mean and HWE scale.
\href{https://github.com/sriramlab/GENIE/blob/22bee9cd53d61eabb2ffcec70acf3c5bb43ede01/src/ge_flexible.cpp#L654-L665}{\texttt{ge\_flexible.cpp}, lines 654--665}.

\bibitem{geniecodetrace}
GENIE source, commit \texttt{22bee9c}, additive trace initialization.
\href{https://github.com/sriramlab/GENIE/blob/22bee9cd53d61eabb2ffcec70acf3c5bb43ede01/src/ge_flexible.cpp#L1655-L1674}{\texttt{ge\_flexible.cpp}, lines 1655--1674}.

\bibitem{geniecodepheno}
GENIE source, commit \texttt{22bee9c}, projected phenotype normalization.
\href{https://github.com/sriramlab/GENIE/blob/22bee9cd53d61eabb2ffcec70acf3c5bb43ede01/src/ge_flexible.cpp#L2331-L2362}{\texttt{ge\_flexible.cpp}, lines 2331--2362}.
\end{thebibliography}

\end{document}
